\documentclass[11pt,a4paper]{article}
% Common preamble for the hint sheets (\input from each file)
\usepackage[utf8]{inputenc}
\usepackage[T1]{fontenc}
\usepackage{lmodern}
\usepackage{amsmath,amssymb,amsthm,mathtools}
\usepackage[a4paper,margin=2.3cm]{geometry}
\usepackage{enumitem}
\usepackage{xcolor}
\usepackage[hidelinks]{hyperref}
\renewcommand{\leq}{\leqslant}
\renewcommand{\le}{\leqslant}
\renewcommand{\geq}{\geqslant}
\renewcommand{\ge}{\geqslant}
\setlength{\parindent}{0pt}
\setlength{\parskip}{0.35em}
% \hintstitle{Sheet title}{Class work / Homework}
\newcommand{\hintstitle}[2]{\begin{center}{\LARGE\bfseries Hints: #1\par}\vspace{0.5em}{\large #2\par}\end{center}
\noindent\textit{The hints to each problem are cumulative. Read one, try again, and open the next only if you are still stuck.}\par\medskip\hrule\medskip}
\newcommand{\problem}[1]{\par\bigskip{\large\bfseries Problem #1}\par\nopagebreak\smallskip}
\newcommand{\subpart}[1]{\par\smallskip\textbf{(#1)}\par\nopagebreak}
\newcommand{\hint}[1]{\par\smallskip\noindent\textcolor{black!70}{\textbf{Hint #1.}}\enspace}
\newcommand{\note}[1]{\par\smallskip\noindent\textit{Note.} #1\par}
\providecommand{\lt}{<}
\providecommand{\gt}{>}
\providecommand{\tightlist}{\setlength{\itemsep}{0pt}\setlength{\parskip}{0pt}}
\emergencystretch=3em

\begin{document}
\hintstitle{Linear algebra -- 4. Functions of matrices and matrix equations}{Class work}

\problem{1}
\hint{1}
Kronecker route: \(X\mapsto AX-XB\) is a linear operator on \(M_n\). Under \(\operatorname{vec}\) its matrix is \(E\otimes A-B^T\otimes E\). Triangularise \(A\) and \(B\) (hence \(B^T\)) simultaneously in the two tensor factors to see that its eigenvalues are \(\lambda_i-\mu_j\). Unique solvability for every \(C\) means this operator is invertible.
\hint{2}
Integral route: let the contour \(\Gamma\) surround \(\operatorname{Spec}A\) but not \(\operatorname{Spec}B\). Use \(A(zE-A)^{-1}=z(zE-A)^{-1}-E\) and the analogous identity for \(B\) to compute \(AX-XB\) for the given integral. Two integrals remain: \(\frac1{2\pi i}\oint(zE-A)^{-1}dz=E\) and \(\oint(zE-B)^{-1}dz=O\).
\hint{3}
Conversely, if \(\lambda\) is a common eigenvalue, \(X=uw^T\) with \(Au=\lambda u\), \(w^TB=\lambda w^T\) is a non-zero solution of \(AX-XB=O\). For the last claim, take \(X=E_{ij}\) (matrix units) in \(AX=XA\).

\problem{2}
\hint{1}
The two matrices must not commute. Try nilpotent ones, whose exponentials are trivial to compute: \(A=E_{12}\), \(B=E_{21}\).
\hint{2}
\(e^Ae^B=(E+E_{12})(E+E_{21})\) has integer entries. \(A+B=\begin{pmatrix}0&1\\1&0\end{pmatrix}\) squares to \(E\), so \(e^{A+B}=\cosh1\,E+\sinh1\,(A+B)\).

\problem{3}
\hint{1}
The target is not diagonalisable. If \(A\) had two distinct eigenvalues, \(A\) and hence \(\sin A\) would be diagonalisable. If \(A\) is scalar, so is \(\sin A\).
\hint{2}
So \(A=\lambda E+N\) with \(N\ne O\), \(N^2=O\), and \(\sin A=\sin\lambda\,E+\cos\lambda\,N\). The eigenvalue condition gives \(\sin\lambda=1\), hence \(\cos\lambda=0\) and \(\sin A=E\). Answer: no such \(A\).

\problem{4}
\hint{1}
A matrix is nilpotent if and only if all its eigenvalues are \(0\). By the spectral mapping principle, the eigenvalues of \(p(e^{AB})\) are \(p(e^{\lambda})\) over the eigenvalues \(\lambda\) of \(AB\).
\hint{2}
\(AB\) and \(BA\) have the same characteristic polynomial, so the same eigenvalues. The statement is true.

\problem{5}
\hint{1}
The left side is \(L(X)=(I+\mathcal N)(X)\), where \(\mathcal N(X)=AX+XA^2\) is the sum of two commuting nilpotent operators on \(M_n(\mathbb R)\) (left multiplication by \(A\) and right multiplication by \(A^2\)). So \(\mathcal N\) is nilpotent and \(L\) is invertible.
\hint{2}
\(X=L^{-1}(A)=\sum_k(-1)^k\mathcal N^k(A)\). Compute \(\mathcal N(A)=A^2\) and \(\mathcal N^2(A)=0\). Answer: \(X=A-A^2\).

\problem{6}
\hint{1}
Write the left side as a linear operator on matrices: \(L=\alpha_{00}I+\sum_{(i,j)\ne(0,0)}\alpha_{ij}\,\mathcal L_A^i\mathcal R_A^j\), where \(\mathcal L_A X=AX\) and \(\mathcal R_AX=XA\).
\hint{2}
\(\mathcal L_A\) and \(\mathcal R_A\) commute and are nilpotent, so the second sum is nilpotent. A non-zero scalar plus a nilpotent operator is invertible.

\problem{7}
\hint{1}
Any solution commutes with \(A=X^2\). First find one obvious square root that is a polynomial in \(A\): with \(M=A-\frac14E\) (so \(M^2=O\)), try \(F=\frac12E+M\).
\hint{2}
\(F=A+\frac14E\) satisfies \(F^2=A\) and is invertible. For any solution \(X\), put \(S=XF^{-1}\). Since \(X\) and \(F\) commute, \(S^2=X^2F^{-2}=E\), and \(S\) commutes with \(A\).
\hint{3}
Conversely, every involution \(S\) commuting with \(A\) gives a solution \(X=S(A+\frac14E)\). Describe these \(S\): either \(S=\pm E\), or \(S\) splits \(\mathbb C^4\) into two \(A\)-invariant \(2\)-dimensional subspaces on which \(A\) is a single Jordan block. In the basis of the sheet, a matrix commuting with \(A\) has \(2\times2\) blocks \(aE+bN\), \(N=E_{12}\).
\end{document}
